\chapter*{Liste des Acronymes}
\addcontentsline{toc}{chapter}{Liste des Acronymes}

% Style pour ce chapitre
\thispagestyle{fancy}
\pagestyle{fancy}

% En-tête stylisé
\begin{center}
\vspace*{1cm}
{\Large\bfseries Liste des Acronymes}\\[0.5cm]
{\small Les acronymes utilisés dans ce document sont définis ci-dessous par ordre alphabétique.}\\[1cm]
\end{center}

% Ligne décorative
\noindent\rule{\textwidth}{1pt}\\[0.5cm]

% Tableau avec design amélioré
\begin{longtable}{p{3.5cm}>{\raggedright\arraybackslash}p{10.5cm}}
\rowcolor{gray!10} % Fond gris clair pour l'en-tête
\textbf{\large Acronyme} & \textbf{\large Signification et Description} \\
\endfirsthead

\rowcolor{gray!10}
\textbf{\large Acronyme} & \textbf{\large Signification et Description} \\
\endhead

\hline
\endfoot

\hline
\multicolumn{2}{c}{\small \textit{Suite de la liste des acronymes à la page suivante}} \\
\endlastfoot

% Contenu du tableau avec meilleure présentation
\textbf{2TUP} & \textbf{2 Tracks Unified Process} (Processus Unifié à Deux Voies) - Méthodologie de développement séparant les aspects fonctionnels et techniques \\
\hline
\textbf{API} & \textbf{Application Programming Interface} (Interface de Programmation d'Application) - Ensemble de règles permettant à des applications de communiquer entre elles \\
\hline
\textbf{APIGEE} & \textbf{API Gateway Enterprise Edition} - Plateforme de gestion d'API développée par Google \\
\hline
\textbf{AWS} & \textbf{Amazon Web Services} - Plateforme de cloud computing d'Amazon \\
\hline
\textbf{BaaS} & \textbf{Backend-as-a-Service} (Backend en tant que Service) - Modèle de cloud computing externalisant les services backend \\
\hline
\textbf{CI/CD} & \textbf{Continuous Integration and Continuous Delivery/Deployment} (Intégration et Livraison/Déploiement Continus) - Pratiques de développement logiciel \\
\hline
\textbf{DAO} & \textbf{Data Access Object} (Objet d'Accès aux Données) - Patron de conception pour l'abstraction de l'accès aux données \\
\hline
\textbf{DTO} & \textbf{Data Transfer Object} (Objet de Transfert de Données) - Patron de conception pour le transfert de données entre les couches \\
\hline
\textbf{FaaS} & \textbf{Function-as-a-Service} (Fonction en tant que Service) - Modèle de cloud computing pour l'exécution de fonctions sans serveur \\
\hline
\textbf{IaC} & \textbf{Infrastructure as Code} (Infrastructure en tant que Code) - Gestion de l'infrastructure via des fichiers de configuration \\
\hline
\textbf{IDMS} & \textbf{Identity Management System} (Système de Gestion d'Identités) - Système pour la gestion des identités utilisateurs \\
\hline
\textbf{JWT} & \textbf{Json Web Token} - Standard pour la création de jetons d'authentification \\
\hline
\textbf{MVC} & \textbf{Model-View-Controller} (Modèle-Vue-Contrôleur) - Patron architectural pour les interfaces utilisateur \\
\hline
\textbf{MVP} & \textbf{Model-View-Presenter} (Modèle-Vue-Présentateur) - Patron architectural dérivé du MVC \\
\hline
\textbf{RUP} & \textbf{Rational Unified Process} - Processus de développement logiciel itératif \\
\hline
\textbf{SDK} & \textbf{Software Development Kit} (Kit de Développement Logiciel) - Ensemble d'outils pour développer des applications \\
\hline
\textbf{SCRUM} & \textbf{Méthodologie Agile de gestion de projet} - Cadre de travail itératif et incrémental \\
\hline
\textbf{SSL} & \textbf{Secure Sockets Layer} - Protocole de sécurisation des communications sur Internet \\
\hline
\textbf{UP} & \textbf{Unified Process} - Processus de développement logiciel unifié \\
\hline
\textbf{VPC} & \textbf{Virtual Private Cloud} (Cloud Privé Virtuel) - Réseau virtuel isolé dans le cloud \\
\hline
\textbf{XP} & \textbf{eXtreme Programming} - Méthodologie de développement logiciel agile \\

\end{longtable}

% Note en bas de page
\vspace{0.5cm}
\begin{center}
\fbox{\parbox{0.9\textwidth}{
\small\textit{Note : Les acronymes sont présentés par ordre alphabétique pour faciliter la recherche. Chaque entrée inclut l'acronyme en gras, sa signification complète et une brève description de son utilisation dans le contexte de ce projet.}
}}
\end{center}